% Options for packages loaded elsewhere
\PassOptionsToPackage{unicode}{hyperref}
\PassOptionsToPackage{hyphens}{url}
\documentclass[
  11pt,
  a4paper,
]{article}
\usepackage{xcolor}
\usepackage[margin=1in]{geometry}
\usepackage{amsmath,amssymb}
\setcounter{secnumdepth}{-\maxdimen} % remove section numbering
\usepackage{iftex}
\ifPDFTeX
  \usepackage[T1]{fontenc}
  \usepackage[utf8]{inputenc}
  \usepackage{textcomp} % provide euro and other symbols
\else % if luatex or xetex
  \usepackage{unicode-math} % this also loads fontspec
  \defaultfontfeatures{Scale=MatchLowercase}
  \defaultfontfeatures[\rmfamily]{Ligatures=TeX,Scale=1}
\fi
\usepackage{lmodern}
\ifPDFTeX\else
  % xetex/luatex font selection
  \setmainfont[]{Latin Modern Roman}
\fi
% Use upquote if available, for straight quotes in verbatim environments
\IfFileExists{upquote.sty}{\usepackage{upquote}}{}
\IfFileExists{microtype.sty}{% use microtype if available
  \usepackage[]{microtype}
  \UseMicrotypeSet[protrusion]{basicmath} % disable protrusion for tt fonts
}{}
\makeatletter
\@ifundefined{KOMAClassName}{% if non-KOMA class
  \IfFileExists{parskip.sty}{%
    \usepackage{parskip}
  }{% else
    \setlength{\parindent}{0pt}
    \setlength{\parskip}{6pt plus 2pt minus 1pt}}
}{% if KOMA class
  \KOMAoptions{parskip=half}}
\makeatother
\usepackage{longtable,booktabs,array}
\newcounter{none} % for unnumbered tables
\usepackage{calc} % for calculating minipage widths
% Correct order of tables after \paragraph or \subparagraph
\usepackage{etoolbox}
\makeatletter
\patchcmd\longtable{\par}{\if@noskipsec\mbox{}\fi\par}{}{}
\makeatother
% Allow footnotes in longtable head/foot
\IfFileExists{footnotehyper.sty}{\usepackage{footnotehyper}}{\usepackage{footnote}}
\makesavenoteenv{longtable}
\setlength{\emergencystretch}{3em} % prevent overfull lines
\providecommand{\tightlist}{%
  \setlength{\itemsep}{0pt}\setlength{\parskip}{0pt}}
\usepackage{mathtools}
\usepackage{enumitem}
\usepackage{microtype}
\usepackage{booktabs}
\usepackage{array}
\usepackage{bookmark}
\IfFileExists{xurl.sty}{\usepackage{xurl}}{} % add URL line breaks if available
\urlstyle{same}
\hypersetup{
  pdftitle={Black Holes in Event Density A Substrate-Level Architecture},
  pdfauthor={Allen Proxmire},
  hidelinks,
  pdfcreator={LaTeX via pandoc}}

\title{Black Holes in Event Density\\
\large A Substrate-Level Architecture}
\author{Allen Proxmire}
\date{May 2026}

\begin{document}
\maketitle

\subsection{Abstract}\label{abstract}

The Event Density (ED) program presents a substrate ontology under which
black-hole physics emerges as a structural consequence of substrate
participation-channel dynamics rather than as a property of a
fundamental spacetime metric. This paper integrates the seven-memo Arc
BH closure (2026-05-01) into a single architectural account of the
black-hole sector. \textbf{Horizons are decoupling surfaces:} under the
Diffusion Coarse-Graining Theorem (DCGT), a horizon is the surface where
cross-participation bandwidth \(\Gamma_\mathrm{cross}\) falls below
hydrodynamic-window resolution, governed by the substrate condition
\(|\nabla\rho|\ell_P^2/\rho_\mathrm{local} \gtrsim \log(R_\mathrm{cg}/\ell_P)\).
The mechanism universalizes across black-hole, Rindler, cosmological,
and acoustic horizons. \textbf{Singularities do not exist:} three
independent substrate constraints --- finite participation (Theorem N1),
discrete micro-events, gradient saturation --- jointly forbid divergent
curvature; the singularity is replaced by a finite-thickness saturated
participation zone. \textbf{Information cannot cross the horizon} as a
participation-bandwidth prohibition (not a causal one); entanglement
amplitude can straddle the surface; \textbf{evaporation is participation
re-routing} through pair-creation-class events, not information
destruction. The four assumptions that generate the standard information
paradox are not imposed at the substrate level, so no paradox arises.
\textbf{Area-law entropy} \(S = (A/\ell_P^2)\log g\) is FORCED in form
by two-dimensional decoupling-surface support + gradient saturation; the
Bekenstein-Hawking \(1/4\) coefficient is INHERITED from substrate motif
counting and is not closed-form derived at present.
\textbf{Wave-black-hole scattering} produces a global phase shift
\(\delta_{\ell m} = \int_\mathrm{path}\Delta\mathcal{A}\,d\lambda\) as a
substrate-derived invariant of minimal ED-channels, with helicity
preservation/flip from anisotropy-basis transport and Kerr twist from
integrated frame-dragging vorticity. The verdict is
\textbf{conditional-positive structural closure} at the architectural
level, parallel to the program's NS-Smoothness Intermediate Path C,
Yang-Mills Clay-relevance, and Arc ED-10 verdicts. Form FORCED / value
INHERITED / extensions OPEN --- the standard methodology of the ED
program --- applied honestly to the black-hole sector.

\begin{center}\rule{0.5\linewidth}{0.5pt}\end{center}

\subsection{1. Introduction}\label{introduction}

\subsubsection{1.1 Motivation}\label{motivation}

Black-hole physics in standard treatments combines general-relativistic
geometry (event horizons, singularities, Killing structure) with
quantum-field-theoretic content (Hawking radiation, Bekenstein-Hawking
entropy, information loss). The combination generates a well-known
cluster of paradoxes: the information-loss paradox (apparent unitarity
violation in evaporation), the firewall puzzle (monogamy-of-entanglement
at the horizon), the complementarity contradiction (interior-vs-exterior
descriptions), and the singularity problem (geodesic incompleteness +
diverging curvature invariants). These puzzles have driven decades of
work on holography, ER=EPR, soft-hair conjectures, fuzzballs, and other
proposals --- none of which has produced a consensus resolution.

The Event Density (ED) program approaches the black-hole sector from a
different ontological starting point. ED's substrate primitives ---
local participation density \(\rho\), participation gradients
\(\nabla\rho\), anisotropy structure, V1 finite-width kernel, P11
commitment irreversibility --- generate spacetime geometry only as a
coarse-grained kinematic reading via the acoustic metric (Arc ED-10),
not as a fundamental object. The substrate ontology does not impose
global unitarity, global Cauchy data, sharp geometric boundaries, or
monogamy-at-boundaries as primitive features. The standard paradoxes are
therefore \textbf{not generated in framework}; they arise only when
continuum-level assumptions are imposed on a fundamentally
local-and-irreversible substrate.

\subsubsection{1.2 ED Ontology in Brief}\label{ed-ontology-in-brief}

ED treats the world as a substrate of participation events. Local
participation density \(\rho(\mathbf{x})\) is the density of viable
participation-channel slots per coarse-graining cell. Participation
events commit irreversibly (Primitive P11): once a substrate channel
commits, the commitment is preserved as part of the local history.
Cross-participation bandwidth \(\Gamma_\mathrm{cross}\) between adjacent
regions controls whether committed structure on one side can register as
committed structure on the other. The hydrodynamic window
\(\ell_P \ll R_\mathrm{cg} \ll L_\mathrm{flow}\) separates substrate
scales from continuum dynamics; the Diffusion Coarse-Graining Theorem
(DCGT, Arc D, 2026-04-30) is the substrate-to-continuum bridge.

\subsubsection{1.3 Why ED Does Not Generate the Standard
Paradoxes}\label{why-ed-does-not-generate-the-standard-paradoxes}

\begin{itemize}
\tightlist
\item
  \textbf{No global unitarity requirement.} Unitarity is a
  coarse-grained continuum property of closed sub-systems in the
  QM-emergence sector (Phase-1, T1--T16); it is not a global constraint
  at the substrate level.
\item
  \textbf{No global Cauchy data.} ED commits events locally; there is no
  spacelike surface on which the entire universe's state must be
  specified.
\item
  \textbf{No sharp geometric boundary.} A horizon in ED is a
  coarse-grained statistical feature of \(\Gamma_\mathrm{cross}\), not a
  knife-edge geometric locus.
\item
  \textbf{No monogamy-at-boundary requirement.} Without a sharp
  boundary, monogamy-of-entanglement enforcement at one is not a
  structural requirement.
\end{itemize}

Removing these four assumptions removes the four standard paradoxes.
\textbf{The information paradox, firewall paradox, complementarity
contradiction, and singularity problem are not solved in ED --- they
simply never arise in its ontology.} What remains is the constructive
question: what does ED's substrate machinery actually deliver for
black-hole architecture?

\subsubsection{1.4 Arc BH as a Consolidation + Derivation
Arc}\label{arc-bh-as-a-consolidation-derivation-arc}

Arc BH (2026-05-01) is a seven-memo arc that consolidates the
interpretive content of ED-06 (information and horizons), ED-10
(curvature emergence), ED-I-06 (fields and forces), and ED-I-26 (wave-BH
scattering channels) into a single substrate-derivation chain grounded
in DCGT + Arc SG (substrate gravity) + Arc ED-10 (acoustic-metric
scalar-tensor covariantization). The seven memos: BH-1 opening; BH-2
horizon-as-decoupling-surface; BH-3 singularity replacement +
strong-curvature audit; BH-4 information vs entanglement + evaporation;
BH-5 area-law entropy + 1/4 coefficient; BH-6 wave-BH scattering; BH-7
synthesis.

\subsubsection{1.5 Summary of Results}\label{summary-of-results}

A single substrate condition unifies the architecture: \[
|\nabla\rho|\ell_P^2/\rho_\mathrm{local} \gtrsim \log(R_\mathrm{cg}/\ell_P).
\] This condition governs horizon formation (BH-2), interior saturation
(BH-3), information blocking (BH-4), participation-capacity saturation
(BH-5), and the strong-curvature scattering region (BH-6). One condition
supports six derivations. The architecture is internally consistent. The
verdict is conditional-positive structural closure: form derived from
substrate machinery; value-layer (entropy coefficient, BHPT spectral
details, \(\beta_\mathrm{crit}\), \(\log g\)) inherited; explicit open
extensions (Hawking spectrum, superradiance amplitude, closed-form
\(\log g\), full Kerr interior) named.

\begin{center}\rule{0.5\linewidth}{0.5pt}\end{center}

\subsection{2. Substrate Primitives and DCGT
Machinery}\label{substrate-primitives-and-dcgt-machinery}

\subsubsection{2.1 Participation Density and Its
Gradients}\label{participation-density-and-its-gradients}

The substrate's load-bearing scalar field is the local participation
density \(\rho(\mathbf{x})\), defined as the density of viable
participation-channel slots per coarse-graining cell of size
\(R_\mathrm{cg}\). Its gradient \(\nabla\rho\) encodes the spatial
variation of the substrate's participation-channel availability. The
dimensionless ratio \[
\sigma(\mathbf{x}) \equiv |\nabla\rho|\ell_P^2/\rho_\mathrm{local}
\] measures the relative steepness of \(\rho\) at the substrate Planck
scale \(\ell_P\). This is the load-bearing quantity for Arc BH:
\(\sigma\) controls horizon formation, interior saturation, and the
strong-curvature scattering region.

\subsubsection{2.2 Anisotropy and Polarization
Structure}\label{anisotropy-and-polarization-structure}

Substrate channels carry directional content: the local anisotropy
tensor encodes the relative weight of participation along different
axes. For minimal ED-channels (gravitational waves, BH-6), the
anisotropy tensor is small and traces out a polarization frame whose
transport along propagation trajectories carries the channel's helicity
content.

\subsubsection{2.3 V1 Finite-Width Kernel}\label{v1-finite-width-kernel}

The V1 vacuum kernel, established in Theorem N1 (2026 foundational
theorem), is the finite-width temporal smearing kernel that mediates
substrate participation events. Its width sets the temporal resolution
of commitment events at a patch and is load-bearing for: (i) the
kernel-arrow result (Theorem 18), (ii) the per-patch motif alphabet
\(g\) in BH-5, (iii) the cross-patch correlation structure that links
adjacent decoupling-surface patches.

\subsubsection{2.4 Multi-Scale Expansion and the Hydrodynamic
Window}\label{multi-scale-expansion-and-the-hydrodynamic-window}

DCGT (Arc D, 2026-04-30) establishes a hydrodynamic-window scale
separation \[
\ell_P \ll R_\mathrm{cg} \ll L_\mathrm{flow}
\] under which substrate dynamics admit a multi-scale expansion to
coarse-grained continuum equations. In Arc BH, DCGT supplies: the
cross-bandwidth \(\Gamma_\mathrm{cross}\) structure, the per-patch
independence assumption used in BH-5 (with corrections from V1
cross-patch correlations), and the channel-propagation kernels used in
BH-6.

\subsubsection{2.5 The Critical Gradient
Condition}\label{the-critical-gradient-condition}

The substrate condition \[
\sigma = |\nabla\rho|\ell_P^2/\rho_\mathrm{local} \gtrsim \log(R_\mathrm{cg}/\ell_P)
\] is the single load-bearing inequality of Arc BH. Below the threshold,
cross-bandwidth between adjacent regions remains finite and substrate
dynamics admit ordinary hydrodynamic-window coarse-graining. At and
above the threshold, \(\Gamma_\mathrm{cross}\) falls below
hydrodynamic-window resolution and the region exhibits
decoupling-surface behavior (Section 3) or gradient-saturation behavior
(Section 4) depending on whether the steep-gradient region is a surface
or a bulk volume.

The threshold value \(\log(R_\mathrm{cg}/\ell_P)\) is itself
substrate-determined (depends on the choice of coarse-graining scale
relative to Planck length); the precise prefactor enters as INHERITED
data labelled \(\beta_\mathrm{crit}\) in subsequent sections.

\begin{center}\rule{0.5\linewidth}{0.5pt}\end{center}

\subsection{3. Horizons as Decoupling Surfaces
(BH-2)}\label{horizons-as-decoupling-surfaces-bh-2}

\subsubsection{3.1 Bandwidth Collapse
Mechanism}\label{bandwidth-collapse-mechanism}

Cross-participation bandwidth between adjacent substrate regions takes
the form (DCGT-derived) \[
\Gamma_\mathrm{cross}(\mathbf{x}) \sim \exp\!\left[-\alpha\,\sigma(\mathbf{x})\right]
\] with \(\alpha\) a substrate-determined dimensionless prefactor. The
exponential suppression is a generic consequence of the multi-scale
expansion: large \(|\nabla\rho|\) at the substrate scale creates steep
barriers to participation-channel mediation between adjacent regions.

\subsubsection{3.2 Decoupling Surface
Definition}\label{decoupling-surface-definition}

The decoupling surface is the locus \[
\Sigma_H \equiv \{\mathbf{x}: \sigma(\mathbf{x}) \gtrsim \log(R_\mathrm{cg}/\ell_P)\},
\] beyond which \(\Gamma_\mathrm{cross}\) has dropped below
hydrodynamic-window resolution. From a coarse-grained continuum
perspective, the substrate channels carrying committed structure across
\(\Sigma_H\) have effectively zero bandwidth.

\subsubsection{3.3 Universalization}\label{universalization}

The same mechanism applies to:

\begin{itemize}
\tightlist
\item
  \textbf{Black-hole event horizons:} steep \(|\nabla\rho|\) from the
  SG-4 / Arc SG participation profile near a high-mass concentration.
\item
  \textbf{Rindler horizons:} acceleration-induced steep \(|\nabla\rho|\)
  in the accelerated frame.
\item
  \textbf{Cosmological horizons:} \(H_0\)-scale gradient set by
  cosmic-horizon participation density.
\item
  \textbf{Acoustic horizons (analog gravity):} flow-induced gradient in
  fluid analog systems.
\end{itemize}

In each case the horizon is a continuum shadow of the same substrate
decoupling surface. This explains, at the substrate level, why these
four classes share thermodynamic-style features (temperature, entropy,
evaporation analogs) despite differing in their continuum-geometry
origins.

\subsubsection{3.4 Geometric Horizon as Continuum
Shadow}\label{geometric-horizon-as-continuum-shadow}

The geometric horizon --- defined in the acoustic-metric reading of Arc
ED-10 as the locus of null-geodesic capture --- is the coarse-grained
continuum signature of \(\Sigma_H\). The two coincide to leading order
in the hydrodynamic expansion; corrections arise at the substrate scale.

\subsubsection{3.5 Observer-Independence}\label{observer-independence}

\(\Sigma_H\) is defined as a level set of the substrate scalar
\(\sigma\), which is built from \(\rho\) and \(\nabla\rho\). Both are
coarse-grained substrate quantities; they do not depend on a choice of
continuum-observer. The decoupling surface is therefore
observer-independent in the same sense that \(\rho\) is.
(Continuum-observer-dependent statements about which side of
\(\Sigma_H\) a particular continuum trajectory is on do depend on the
observer, as expected --- but the surface itself does not.)

\begin{center}\rule{0.5\linewidth}{0.5pt}\end{center}

\subsection{4. No Singularities: Interior Saturation
(BH-3)}\label{no-singularities-interior-saturation-bh-3}

\subsubsection{4.1 Three Substrate
Constraints}\label{three-substrate-constraints}

Three independent substrate constraints jointly forbid divergent
curvature/energy/tidal/incompleteness at the substrate level:

\textbf{(C1) Finite participation (P4 + Theorem N1).} The substrate has
finite participation-channel density per coarse-graining cell. There is
no substrate state in which infinitely many channels participate at a
single point.

\textbf{(C2) Discrete micro-events (P01).} Substrate events are
discrete: each commitment is a finite-amplitude occurrence at a
finite-resolution patch. Continuum singularity statements that require
unbounded micro-event multiplicity are inadmissible.

\textbf{(C3) Gradient saturation (P4 mobility-capacity bound).}
\(|\nabla\rho|\) has a substrate-imposed maximum: when
\(\sigma \to \beta_\mathrm{crit}\), the substrate's mobility-capacity
bound caps further steepening. Gradients cannot run away to infinity.

Any one of (C1)--(C3) alone would suffice to forbid a continuum
singularity; the three together overdetermine the result.

\subsubsection{4.2 Acoustic Metric
Breakdown}\label{acoustic-metric-breakdown}

The acoustic-metric scalar-tensor reading of Arc ED-10 is a
coarse-grained kinematic structure:
\(g_{\mu\nu}^\mathrm{ac}[\rho,\nabla\rho]\). Its validity requires the
hydrodynamic-window expansion of DCGT to be controlled. When
\(\sigma \gtrsim \beta_\mathrm{crit}\), the multi-scale expansion's
small parameter ceases to be small; the acoustic metric ceases to be a
controlled approximation. \textbf{The acoustic metric breaks down before
the substrate does.}

This is the substrate-level meaning of ``singularity'': the breakdown is
a property of the kinematic continuum reading, not of the underlying
substrate. Substrate dynamics remain finite throughout.

\subsubsection{4.3 Replacement: Finite-Thickness Saturated
Zone}\label{replacement-finite-thickness-saturated-zone}

The substrate replacement for the singular endpoint is a
finite-thickness saturated participation zone characterized by:

\begin{itemize}
\tightlist
\item
  Gradients at the substrate-imposed maximum:
  \(\sigma \to \beta_\mathrm{crit}\).
\item
  Cross-participation bandwidth collapsed in all directions (not just
  the radial-equivalent of \(\Sigma_H\)).
\item
  Only the commitment-order direction (P11) remaining structurally
  meaningful.
\item
  Acoustic metric degenerate or undefined.
\item
  Substrate dynamics finite, consisting of saturated commitment cycles.
\end{itemize}

\subsubsection{4.4 No Geodesic Incompleteness at Substrate
Level}\label{no-geodesic-incompleteness-at-substrate-level}

Geodesic incompleteness is a property of the acoustic-metric reading. At
the substrate level there are no geodesics --- there are only
commitment-event sequences. Commitment-event sequences inside the
saturated zone proceed in finite-amplitude finite-resolution increments;
there is no sequence that ``ends'' at a singular point. Substrate-level
evolution is everywhere defined.

\begin{center}\rule{0.5\linewidth}{0.5pt}\end{center}

\subsection{5. Information Architecture and Evaporation
(BH-4)}\label{information-architecture-and-evaporation-bh-4}

\subsubsection{5.1 Committed vs Uncommitted
Structure}\label{committed-vs-uncommitted-structure}

ED distinguishes two structurally different forms of correlation:

\begin{itemize}
\tightlist
\item
  \textbf{Committed structure (information):} P11-irreversible records
  of participation events. Transmission of committed structure between
  two regions requires nonzero \(\Gamma_\mathrm{cross}\) between them.
\item
  \textbf{Uncommitted structure (entanglement):} correlation potential
  between two regions, set up at pair-formation. Persistence of
  uncommitted structure does not require ongoing cross-bandwidth.
\end{itemize}

This is the substrate translation of the standard quantum-information
no-signaling-but-correlations-persist behavior.

\subsubsection{5.2 Why Information Cannot Cross a
Horizon}\label{why-information-cannot-cross-a-horizon}

At \(\Sigma_H\), \(\Gamma_\mathrm{cross}\) is below hydrodynamic-window
resolution (Section 3). Therefore: - Committed structure transmission
across \(\Sigma_H\) requires bandwidth that is structurally absent. -
Information, in the substrate sense, does not pass through a horizon.

This is \textbf{not a causal prohibition} (ED has no fundamental
light-cone primitive); it is a \textbf{participation-bandwidth
prohibition}.

\subsubsection{5.3 Entanglement
Straddling}\label{entanglement-straddling}

Uncommitted structure does not require ongoing
\(\Gamma_\mathrm{cross}\); it is set by joint participation amplitude at
pair-formation. Pair-amplitude structures can therefore exist with one
leg inside and one leg outside \(\Sigma_H\) without violating Section
5.2.

\subsubsection{5.4 Evaporation as Participation
Re-Routing}\label{evaporation-as-participation-re-routing}

The evaporation mechanism follows from the architecture above:

\begin{enumerate}
\def\labelenumi{\arabic{enumi}.}
\tightlist
\item
  Horizon blocks committed structure.
\item
  Uncommitted pair structures can straddle \(\Sigma_H\).
\item
  Strong \(|\nabla\rho|\) near \(\Sigma_H\) supports pair-creation-class
  events (gradient-driven, substrate analog of strong-field vacuum pair
  production).
\item
  Commitment events on these pairs commit one leg inward (joining the
  BH-3 saturated zone) and the other outward (joining the radiation
  field).
\item
  Joint correlation is preserved globally across the inward + outward
  commitments.
\end{enumerate}

This is \textbf{participation re-routing}, not information destruction.
What an external observer reads as thermal Hawking radiation is the
coarse-grained statistics of outward-committed legs of these pairs. (The
thermal spectrum derivation itself is reserved for the next arc ---
priority item B4 --- and is explicitly out of scope for the present
paper.)

\subsubsection{5.5 No Information Paradox}\label{no-information-paradox}

The four assumptions that generate the standard information paradox are:
(a) global unitarity, (b) global Cauchy data, (c) sharp geometric
boundary, (d) monogamy-at-boundary.

None of these is imposed at the substrate level. The paradox does not
arise. The same applies to firewalls and complementarity: both require a
sharp geometric boundary, which the ED ontology does not provide.

The ED stance is structural: BH puzzles are \textbf{not solved in
framework} but \textbf{not generated in framework}. The same pattern
applies to the singularity (BH-3) and to the horizon (BH-2): standard
puzzles dissolve when the continuum-level assumptions that generate them
are not imposed.

\begin{center}\rule{0.5\linewidth}{0.5pt}\end{center}

\subsection{6. Area-Law Entropy and the 1/4 Coefficient
(BH-5)}\label{area-law-entropy-and-the-14-coefficient-bh-5}

\subsubsection{6.1 Participation Capacity}\label{participation-capacity}

Define the participation-capacity density \(\chi(\rho,\nabla\rho)\) as
the substrate-level density of viable commitment-channel slots per
Planck-area patch on \(\Sigma_H\). Total participation capacity: \[
\mathcal{C}_H = \int_{\Sigma_H} \frac{dA}{\ell_P^2}\,\chi(\rho,\nabla\rho).
\] Two structural points: (i) area-extensivity (the integrand is
supported only on \(\Sigma_H\), not on a bulk volume); (ii) saturation
in the BH-relevant regime (\(\chi \to 1\) where gradient saturation
holds, per BH-3).

\subsubsection{6.2 Area-Law Form}\label{area-law-form}

In the BH-relevant regime, \[
\mathcal{C}_H \approx A/\ell_P^2,
\] with \(A\) the coarse-grained acoustic-metric area of \(\Sigma_H\).
The area-law form is \textbf{derived} from substrate machinery:
two-dimensional surface support + gradient saturation.

\subsubsection{6.3 Entropy as Log of Viable Commitment
Histories}\label{entropy-as-log-of-viable-commitment-histories}

ED's substrate definition:
\(S = \log(\text{number of viable commitment histories compatible with the macroscopic constraint})\).
With per-patch finite alphabet \(g\) (the substrate's local count of
distinct commitment motifs per Planck-area patch) and per-patch
independence under DCGT: \[
S = \mathcal{C}_H \log g \approx \frac{A}{\ell_P^2}\log g.
\]

\subsubsection{6.4 The 1/4 Coefficient:
Inherited}\label{the-14-coefficient-inherited}

The Bekenstein-Hawking value \(S_\mathrm{BH} = A/(4\ell_P^2)\)
corresponds to \(\log g = 1/4\), i.e.~\(g = e^{1/4} \approx 1.284\).
This is \textbf{not an integer}. Naive per-patch-independent
finite-alphabet counting cannot reproduce \(1/4\): integer \(g \geq 2\)
all give \(\log g \geq \ln 2 \approx 0.693\).

To recover \(1/4\), the per-patch-independence assumption must be
relaxed: adjacent patches share commitment-history correlations through
V1-kernel cross-patch overlap. The effective per-patch contribution is
reduced from \(\log g\) (independent count) to a smaller correlated
value. Recovering \(1/4\) exactly requires substrate-level state
counting that is not currently in hand --- structurally similar in
difficulty to closed-form derivation of \(\kappa/|\hat{N}'| = 0.001766\)
(priority item E4).

\textbf{Verdict on the 1/4 coefficient: INHERITED, not derived.}
Form-FORCED + value-INHERITED methodology applied honestly. Same
category as \(a_0\), \(G\), \(\Lambda\), gauge couplings, and the other
substrate-determined dimensionless numbers in the ED program.

\subsubsection{6.5 Implications for Motif Counting and V1
Correlations}\label{implications-for-motif-counting-and-v1-correlations}

The closed-form \(\log g\) derivation would require either: (a) an
explicit V1-kernel-counting argument analogous to a substrate-level
partition function; or (b) a derivation of the per-patch commitment
alphabet from microscopic substrate dynamics. Either route is a
candidate future arc. Currently flagged as \textbf{OPEN item O2}.

\begin{center}\rule{0.5\linewidth}{0.5pt}\end{center}

\subsection{7. Wave-Black-Hole Scattering
(BH-6)}\label{wave-black-hole-scattering-bh-6}

\subsubsection{7.1 Minimal ED-Channel}\label{minimal-ed-channel}

A minimal ED-channel (gravitational wave, in the continuum reading) is a
substrate excitation characterized by: - Low multiplicity (small number
of participation degrees of freedom involved per coarse-graining cell).
- Low anisotropy (anisotropy tensor magnitude small compared to
substrate-saturation). - DCGT-coherent propagation (channel survives
multi-scale expansion to a wave equation on the acoustic metric).

Channel global phase along propagation trajectory \(\gamma(\lambda)\):
\[
\Phi[\gamma] = \int_\gamma \mathcal{A}(\rho,\nabla\rho,\text{anisotropy})\,d\lambda.
\]

\subsubsection{7.2 Maximal ED-Structure}\label{maximal-ed-structure}

A maximal ED-structure (black hole) is high-multiplicity,
saturated-gradient, decoupling-surface-bounded, with curvature-like
participation structure given by SG-4 / Arc SG. In the BH-relevant
regime, scattered channels propagate through a region where
\(\sigma \gg 1\) near \(\Sigma_H\) and the acoustic metric is strongly
curved.

\subsubsection{7.3 Phase Shift as Global
Invariant}\label{phase-shift-as-global-invariant}

The minimal channel's propagation kernel is distorted by the maximal
structure's saturated gradients. The channel's phase accumulates a shift
relative to its asymptotic free-wave value: \[
\delta_{\ell m} = \int_\mathrm{path}\Delta\mathcal{A}\,d\lambda
\] with
\(\Delta\mathcal{A} = \mathcal{A}_\mathrm{full} - \mathcal{A}_\mathrm{asymptotic}\)
the curvature-induced excess action density along the trajectory. The
shift is a single scalar per partial-wave mode \((\ell,m)\) because the
channel is minimal; it is inevitable along any trajectory threading the
strong-curvature region because the gradients are saturated.

\subsubsection{7.4 S-Matrix
Exponentiation}\label{s-matrix-exponentiation}

The S-matrix structure \(S = e^{iN}\) from ED-I-26 is recovered: \(N\)
is the integrated channel action; \(\delta_{\ell m}\) are its
partial-wave components. Exterior-side unitarity follows from BH-4 (no
committed structure leaks across \(\Sigma_H\), so no probability is lost
from the scattering channel; entanglement-amplitude correlations with
the interior persist but do not enter exterior unitarity bookkeeping).

\subsubsection{7.5 Helicity Preservation and
Flip}\label{helicity-preservation-and-flip}

For axisymmetric, non-rotating backgrounds (Schwarzschild-like):
principal directions of the participation profile are aligned uniformly
along the trajectory's symmetry plane. The channel's anisotropy tensor
returns to itself up to an overall scalar phase under parallel
transport. \textbf{Helicity is preserved}; phase shifts act identically
on both helicity states.

For backgrounds with frame-dragging (Kerr-like): principal directions
rotate along the trajectory. The channel's anisotropy tensor picks up a
twist relative to a non-rotating reference frame: \[
\Delta\phi_\mathrm{twist} = \int \omega_\mathrm{FD}\,d\lambda.
\] \textbf{Helicity flip channels open}, with mixing strength tied to
spin parameter via \(\omega_\mathrm{FD}\).

\subsubsection{7.6 Kerr Twist and
Superradiance}\label{kerr-twist-and-superradiance}

The Kerr phase shift relative to Schwarzschild has a substrate origin:
frame-dragging is a global rotation of the participation anisotropy
basis, sourced by substrate-level vorticity in the participation flow of
a spinning maximal structure. The channel acquires a twist proportional
to integrated background vorticity along its trajectory.
Trajectory-dependence of the integration produces the
prograde/retrograde scattering asymmetry.

\textbf{Superradiance} (energy extraction by certain scattering
channels) is structurally compatible with this picture: channels
traversing the rotating-vorticity exterior can net positive energy
transfer from background to channel via the participation-flow
vorticity. No committed structure is extracted from the interior; the
energy transfer is mediated by exterior participation flow. Full
amplitude derivation is \textbf{OPEN item O1}.

\begin{center}\rule{0.5\linewidth}{0.5pt}\end{center}

\subsection{8. Unified Black-Hole
Architecture}\label{unified-black-hole-architecture}

The six derivations integrate into a single architectural picture
supported by the single substrate condition
\(\sigma \gtrsim \log(R_\mathrm{cg}/\ell_P)\):

\begin{itemize}
\tightlist
\item
  \textbf{Formation.} Steep participation-density gradients drive
  \(\sigma\) across threshold; cross-bandwidth collapses; decoupling
  surface \(\Sigma_H\) emerges (Section 3 / BH-2).
\item
  \textbf{Interior.} Gradients continue steepening but saturate at the
  substrate-imposed maximum (C3); finite-thickness saturated zone
  replaces the singular endpoint (Section 4 / BH-3).
\item
  \textbf{Information flow.} Committed structure cannot cross
  \(\Sigma_H\) (bandwidth-blocked); entanglement amplitude can straddle
  (Section 5 / BH-4).
\item
  \textbf{Thermal behavior.} Asymmetric participation flow at the
  saturated decoupling surface produces, at the coarse-grained continuum
  level, the statistical features standardly read as thermal (mechanism:
  Section 5; spectrum derivation: deferred to B4).
\item
  \textbf{Entropy.} Area-law form from participation capacity of
  \(\Sigma_H\) (Section 6 / BH-5); coefficient INHERITED.
\item
  \textbf{Scattering.} Minimal ED-channel acquires global phase shift
  through saturated-gradient exterior; helicity behavior + Kerr twist
  follow from anisotropy-basis transport (Section 7 / BH-6).
\item
  \textbf{Evaporation.} Pair-creation-class events at \(\Sigma_H\) route
  participation outward + inward; correlation structure preserved
  globally (Section 5 / BH-4).
\end{itemize}

The architecture is internally consistent: one substrate-level condition
governs all six. \textbf{No singularity. No information paradox. No
global unitarity requirement. No firewall. No complementarity
contradiction. No geometric-boundary primitive.}

\begin{center}\rule{0.5\linewidth}{0.5pt}\end{center}

\subsection{9. FORCED / INHERITED / OPEN
Classification}\label{forced-inherited-open-classification}

\subsubsection{FORCED (substrate-level
form)}\label{forced-substrate-level-form}

\begin{itemize}
\tightlist
\item
  Horizon as decoupling surface (BH-2).
\item
  No singularities (BH-3).
\item
  Information blocking + entanglement straddling (BH-4).
\item
  Evaporation as participation re-routing (BH-4).
\item
  Area-law form (BH-5).
\item
  BHPT phase shift structure (BH-6).
\item
  Helicity behavior under axisymmetric vs frame-dragging backgrounds
  (BH-6).
\item
  Kerr twist as integrated background vorticity (BH-6).
\end{itemize}

\subsubsection{INHERITED (value-level)}\label{inherited-value-level}

\begin{itemize}
\tightlist
\item
  1/4 entropy coefficient (BH-5): corresponds to \(\log g\) from
  substrate motif counting; closed-form value not currently extracted.
\item
  BHPT spectral details (BH-6): match to specific \(\delta_{\ell m}\)
  values conditional on Arc ED-10 + Arc SG acoustic-metric
  prerequisites.
\item
  \(\beta_\mathrm{crit}\) (BH-2 / BH-3): substrate-determined
  dimensionless threshold; precise value depends on coarse-graining
  choices.
\item
  V1 kernel width + per-patch motif alphabet \(g\) (BH-5): substrate
  constants whose closed-form values are not currently extracted.
\end{itemize}

\subsubsection{OPEN (future arcs)}\label{open-future-arcs}

\begin{itemize}
\tightlist
\item
  \textbf{B4 --- Hawking spectrum derivation} (next natural arc).
  Mechanism in hand from BH-4; load-bearing work is the explicit V5
  cross-chain correlation calculation that produces the spectrum. Match
  to \(T_H = \kappa/(2\pi)\) would be structural recovery; deviation a
  falsifier.
\item
  \textbf{O1 --- Superradiance amplitude derivation.} Structurally
  compatible per BH-6; full derivation requires extending BH-6 + Arc SG.
\item
  \textbf{O2 --- Closed-form \(\log g\) from V1 kernel + motif
  counting.} Substrate-level state-counting derivation, structurally
  similar in difficulty to E4.
\item
  \textbf{O3 --- Full Kerr interior audit.} Extension of BH-3 to
  rotating-BH interior architecture (ergosphere, Cauchy horizons,
  frame-dragging at saturation).
\end{itemize}

\begin{center}\rule{0.5\linewidth}{0.5pt}\end{center}

\subsection{10. Discussion and Outlook}\label{discussion-and-outlook}

\subsubsection{10.1 Implications for Quantum
Gravity}\label{implications-for-quantum-gravity}

The standard quantum-gravity program seeks to reconcile
general-relativistic geometry with quantum-field-theoretic content
within a single dynamical framework, treating the metric as a
fundamental (or holographically dual) degree of freedom. ED takes a
different starting point: the metric is a coarse-grained kinematic
reading; the substrate is local participation events with irreversible
commitment. From this starting point, the standard black-hole puzzles
are not generated. Whether this dissolution is satisfying or merely
re-locates the difficulty is a judgement call. The honest claim is that
ED delivers a substrate-level architecture for the black-hole sector
that is internally consistent, derives form from substrate machinery,
and inherits value-layer details from substrate constants --- the same
form-FORCED + value-INHERITED pattern that holds across the rest of the
program.

\subsubsection{10.2 Relation to GR, QFT, and Analog
Gravity}\label{relation-to-gr-qft-and-analog-gravity}

\begin{itemize}
\tightlist
\item
  \textbf{GR.} ED does not derive the Einstein equation. The
  acoustic-metric reading is kinematic, not full GR. Black-hole
  architecture in ED is \textbf{not} a derivation of GR black-hole
  solutions; it is a substrate-level account of how black-hole-like
  structure emerges from substrate participation dynamics, with
  continuum-geometry features matching GR-class predictions to leading
  order via the acoustic-metric scalar-tensor approximation (Arc ED-10).
\item
  \textbf{QFT.} ED supplies QM emergence (Phase-1, T1--T16) and
  form-level QFT content (T17, T18) within the canonical-ED hydrodynamic
  window. Hawking-radiation derivation (priority B4) is the natural site
  to test how far the form-level QFT machinery extends into the
  strong-curvature regime.
\item
  \textbf{Analog gravity.} The universalization of the
  decoupling-surface mechanism across BH / Rindler / cosmological /
  acoustic horizons (Section 3) is a substrate-level explanation for the
  empirical similarity of analog-gravity systems to gravitational ones.
  Acoustic horizons are not ``weakly gravity-like''; they are decoupling
  surfaces of the same substrate type.
\end{itemize}

\subsubsection{10.3 Why ED Avoids Paradoxes Rather Than Solving
Them}\label{why-ed-avoids-paradoxes-rather-than-solving-them}

\textbf{The information paradox, firewall paradox, complementarity
contradiction, and singularity problem are not solved in ED --- they
simply never arise in its ontology.} The standard paradox-cluster
(information loss, firewalls, complementarity) requires four
assumptions: global unitarity, global Cauchy data, sharp geometric
boundary, monogamy-at-boundary. None is imposed at the substrate level.
The dissolution is structural, not constructive: ED does not propose a
clever resolution within a paradox-generating framework; it identifies
the assumptions whose joint imposition generates the paradox and shows
that the substrate ontology does not impose them. The same pattern
appears in BH-3 (singularity) and BH-2 (horizon). It is a recurring
feature of the ED program: standard puzzles dissolve when
continuum-level assumptions are not imposed on a fundamentally
local-and-irreversible substrate.

\subsubsection{10.4 Next Arc: Hawking Spectrum
(B4)}\label{next-arc-hawking-spectrum-b4}

The next natural arc is Hawking-spectrum derivation. The mechanism is in
hand from BH-4 (asymmetric participation flow at saturated decoupling
surface produces outward-radiation statistics). The load-bearing
technical work is the explicit V5 cross-chain correlation calculation
that produces the spectrum. A match to \(T_H = \kappa/(2\pi)\) would be
a structural recovery; a deviation would be a falsifier of the
architecture. Estimated 6--8 memos. Couples naturally with O1
(superradiance) since both involve substrate-level cross-chain
correlations at the decoupling surface.

\subsubsection{10.5 Long-Term Program
Trajectory}\label{long-term-program-trajectory}

Arc BH closure brings the structurally-complete column to: QM emergence
(Phase-1) + form-level QFT (T17, T18) + substrate gravity (Arc SG / T19
/ T20 / ECR / T21) + curvature emergence (Arc ED-10) + fluid dynamics
(NS / MHD) + non-Abelian gauge theory (Yang-Mills) + soft-matter
mobility (UDM / P4-NN) + black-hole architecture (Arc BH). Eight sectors
with form-derived content. The remaining major open structural fronts
are: GR-4A (Einstein equation emergence, SPECULATIVE), Arc COSMO (cosmic
expansion / \(H_0\) derivation, SPECULATIVE), and B4 (Hawking spectrum,
the next concrete arc). The program-level methodology --- form-FORCED +
value-INHERITED + explicit OPEN items --- is consistent across every
closed arc.

\begin{center}\rule{0.5\linewidth}{0.5pt}\end{center}

\subsection{11. Appendices}\label{appendices}

\subsubsection{Appendix A --- DCGT Derivation
Details}\label{appendix-a-dcgt-derivation-details}

DCGT (Diffusion Coarse-Graining Theorem, Arc D, 2026-04-30) supplies the
substrate-to-continuum bridge for canonical-ED dynamical content. The
hydrodynamic-window scale separation
\(\ell_P \ll R_\mathrm{cg} \ll L_\mathrm{flow}\) admits a multi-scale
expansion. At leading order, scalar diffusion + directional viscosity +
V1→R1 + V5→Maxwell memory + T17 minimal-coupling Lorentz are all derived
from substrate primitives.

For Arc BH, DCGT supplies: (i) the cross-bandwidth structure
\(\Gamma_\mathrm{cross} \sim \exp[-\alpha\sigma]\) used in Section 3;
(ii) the per-patch independence assumption used in Section 6 (with V1
cross-patch correction terms responsible for reducing naive
integer-\(g\) counts to the inherited \(\log g \approx 1/4\)); (iii) the
channel-propagation kernels used in Section 7. The full derivation is in
Arc D's six memos at \texttt{theory/Arc\_D/}.

\subsubsection{Appendix B --- Participation Capacity and Motif
Counting}\label{appendix-b-participation-capacity-and-motif-counting}

The participation-capacity density \(\chi(\rho,\nabla\rho)\) measures
viable commitment-channel slots per Planck-area patch. By construction
\(\chi \in [0,1]\), with \(\chi \to 1\) in the gradient-saturated
regime. The total participation capacity is the surface integral
\(\mathcal{C}_H = \int_{\Sigma_H}(dA/\ell_P^2)\chi\).

The per-patch motif alphabet \(g\) counts distinct viable commitment
motifs per Planck-area patch on a saturated decoupling surface.
Constraints on \(g\):

\begin{itemize}
\tightlist
\item
  P11 (each commitment slot is a one-way switch): bounds motif count by
  the count of distinguishable commitment-switch sequences.
\item
  V1 finite kernel width: bounds the count of distinguishable temporal
  orderings within a patch.
\item
  Gradient saturation (BH-3): suppresses motif bleeding between adjacent
  patches, justifying per-patch independence to leading order.
\item
  DCGT hydrodynamic-window resolution: distinguishability requires
  motifs to survive coarse-graining over \(R_\mathrm{cg}\).
\end{itemize}

These constraints determine \(g\) as a finite substrate-determined
integer ≥ 2. Closed-form extraction is OPEN item O2.

\subsubsection{Appendix C --- Scattering Integrals and Anisotropy
Transport}\label{appendix-c-scattering-integrals-and-anisotropy-transport}

The phase-shift integral
\(\delta_{\ell m} = \int_\mathrm{path}\Delta\mathcal{A}\,d\lambda\) has
the structure of a path-integrated curvature-coupling correction. The
integrand
\(\Delta\mathcal{A} = \mathcal{A}_\mathrm{full} - \mathcal{A}_\mathrm{asymptotic}\)
is built from the channel's substrate action density evaluated against
the full SG-4 / Arc SG participation profile minus its asymptotic
flat-space limit.

Anisotropy-basis transport along the trajectory is governed by parallel
transport of the channel's polarization tensor through the local
participation-frame. For axisymmetric backgrounds, parallel transport
closes (helicity preserved). For frame-dragging backgrounds, parallel
transport accumulates a twist
\(\Delta\phi_\mathrm{twist} = \int\omega_\mathrm{FD}\,d\lambda\)
proportional to integrated substrate vorticity.

\subsubsection{Appendix D --- Comparison with GR / QFT BH
Structures}\label{appendix-d-comparison-with-gr-qft-bh-structures}

{\def\LTcaptype{none} % do not increment counter
\begin{longtable}[]{@{}
  >{\raggedright\arraybackslash}p{(\linewidth - 4\tabcolsep) * \real{0.3333}}
  >{\raggedright\arraybackslash}p{(\linewidth - 4\tabcolsep) * \real{0.3333}}
  >{\raggedright\arraybackslash}p{(\linewidth - 4\tabcolsep) * \real{0.3333}}@{}}
\toprule\noalign{}
\begin{minipage}[b]{\linewidth}\raggedright
Feature
\end{minipage} & \begin{minipage}[b]{\linewidth}\raggedright
GR/QFT
\end{minipage} & \begin{minipage}[b]{\linewidth}\raggedright
ED (this paper)
\end{minipage} \\
\midrule\noalign{}
\endhead
\bottomrule\noalign{}
\endlastfoot
Horizon & Null hypersurface in metric & Decoupling surface in
\(\sigma\) \\
Singularity & Geodesic incompleteness + diverging curvature &
Acoustic-metric breakdown; substrate finite \\
Information & Crosses horizon classically; paradox quantum & Cannot
cross (bandwidth); paradox not generated \\
Entanglement & Monogamy-at-boundary forces firewall & Straddles surface;
no firewall \\
Entropy & \(A/(4\ell_P^2)\) derived (e.g.~holography) & Form
\(A/\ell_P^2 \cdot \log g\) derived; coefficient inherited \\
Hawking \(T\) & \(\kappa/(2\pi)\) from QFT in curved background &
Mechanism from re-routing; spectrum derivation OPEN (B4) \\
Helicity preservation & Schwarzschild diagonal & Anisotropy-basis
aligned \\
Kerr twist & Frame-dragging in metric & Integrated substrate
vorticity \\
Superradiance & Penrose process / wave amplification & Structurally
compatible; OPEN (O1) \\
Universal across horizon types & Analog gravity heuristically related &
Universal at substrate level (Section 3.3) \\
\end{longtable}
}

\begin{center}\rule{0.5\linewidth}{0.5pt}\end{center}

\subsection{12. Metadata}\label{metadata}

\textbf{Keywords.} Event Density; Black Holes; Substrate Physics;
Decoupling Surface; DCGT; Acoustic Metric; Information Paradox;
Bekenstein-Hawking Entropy; BHPT; Kerr Scattering; Form-FORCED /
Value-INHERITED Methodology.

\textbf{Version.} 1.0 --- May 2026.

\textbf{Arc BH References.} Seven memos in
\texttt{theory/Black\_Holes/}: Arc\_BH\_1\_Opening;
Arc\_BH\_2\_Horizon\_DCGT;
Arc\_BH\_3\_Singularity\_And\_Strong\_Curvature\_Audit;
Arc\_BH\_4\_Information\_And\_Evaporation;
Arc\_BH\_5\_Area\_Law\_Entropy; Arc\_BH\_6\_Wave\_BH\_Scattering;
Arc\_BH\_7\_Synthesis.

\textbf{Related ED Papers.} \emph{Foundations of Substrate Gravity}
(April 2026, two complementary papers); \emph{Event Density Foundations
of Quantum Field Theory: A Unified Substrate-to-Continuum Overview}
(April 2026); \emph{Navier-Stokes Synthesis Paper} (April 2026, with
Appendices C/D/E covering MHD + DCGT + Yang-Mills); \emph{Event Density:
One Substrate, Three Domains} (April 2026); \emph{Theorem 17:
Gauge-Field-as-Rule-Type}; \emph{Theorem 18: V1 Kernel Retardation}.

\textbf{Citation Block (Zenodo-safe).}

\begin{quote}
Proxmire, A. (2026). \emph{Black Holes in Event Density: A
Substrate-Level Architecture.} Event Density Program, Black Hole
Foundations Paper, May 2026.
\end{quote}

\textbf{Conditional-Positive Verdict Disclaimer.} This paper presents a
substrate-level architectural account of black-hole physics.
Form-derivations are FORCED from substrate machinery; value-layer
details (entropy coefficient, BHPT spectral values,
\(\beta_\mathrm{crit}\), \(\log g\)) are INHERITED and not closed-form
derived at present. No claim of GR derivation is made; the
acoustic-metric reading is kinematic per ED-Phys-10 guardrails. Open
extensions (Hawking spectrum derivation, superradiance amplitude,
closed-form \(\log g\), full Kerr interior audit) are explicitly named.

\end{document}
